% ===================================================================
%  TABLEAU N° 12 — La Gare. --- Les Chemins de fer. --- Les Bateaux.
%  Feuillets 71 a 74 du fac-simile = folios 69 a 72.
%  Correspondance : folio imprime = numero de feuillet - 2.
%
%  Les cles « %%K » sont les MEMES que dans texto/io/21-tabelo-12.tex :
%  c'est par elles que la page de lecture met les deux textes en
%  regard. La page d'ouverture (feuillet 71) porte un « NOTA. » en
%  petites capitales suivi d'un texte italique, exactement comme le
%  « Noto. » ido correspondant (feuillet 82 de l'ido) ; il precede
%  l'alinea « 1. » et est donc releve comme matiere continue du bloc
%  %%K t12-apar-1, sans cle propre, a l'image de la note de bas de
%  page du tableau 1 (feuillet 9) dans le modele.
%
%  ATTENTION — LE FAC-SIMILE FRANCAIS EST UNE PRISE DE VUE, contraste
%  inegal d'une page a l'autre ; plusieurs renvois chiffres a demi
%  effaces (store, fourgon, bouillottes, ponton, petite) ont ete
%  recoupes avec la legende numerique de texto/io/21-tabelo-12.tex,
%  qui porte les memes numeros d'objets.
% ===================================================================

% --- Feuillet 71 (folio 69 : ouverture de tableau) ------------------
\begin{VUpage}[71]{69}
%%K t12-apar-1 apar
\VUsaut{4.0mm}
\VUtitre{15.0pt}{60}{TABLEAU N\textsuperscript{o} 12}
\VUsaut{2.4mm}
% Le feuillet ne porte AUCUN fleuron entre les deux lignes de titre :
% la petite tache d'encre visible a gauche de « Les Bateaux. » est un
% accident d'impression, non un ornement. L'interligne mesure au
% facsimile environ 16pt pour un corps de 11.4pt, d'ou le blanc ci-
% dessous. Les deux lignes forment un seul titre, que la page de
% lecture recolle avec un cadratin, comme « La Staciono. --- La
% Fervoyi. --- La Navi. » cote ido.
\VUtitre{11.4pt}{0}{La Gare. --- Les Chemins de fer.}
\VUsaut{0.8mm}
\VUtitre{11.4pt}{0}{Les Bateaux.}
\VUsaut{3.4mm}
\VUblancAlinea
\textsc{Nota.} --- \textit{Les horaires choisis pour chaque langue}\nl
\textit{étant différents, nous avons dû prendre comme}\nl
\textit{exemple l'horaire du tableau} \VUgras{français.}

%%K t12-01-1 p
\VUblancAlinea
1. --- Je viens de faire un voyage en Allemagne.\nl
Avant de l'entreprendre, j'avais acheté un \VUgras{indi}\cc
\VUgras{cateur} \textsuperscript{(15)} pour savoir l'heure de départ des trains.\nl
Ayant constaté que le train le plus commode partait\nl
de Paris à 9 heures du soir et arrivait à Strasbourg\nl
à 1 h. 13, je fis mes préparatifs de voyage et, le lende\cc
main, je me fis conduire à la gare.

%%K t12-02-1 p
\VUblancAlinea
2. --- Lorsque je pénétrai dans la grande salle de\nl
départ, derrière un vieux \VUgras{prêtre} \textsuperscript{(2)} de campagne\nl
qui portait à la main son \VUgras{sac de nuit} \textsuperscript{(3)}, un grand\nl
nombre de \VUgras{voyageurs} \textsuperscript{(1)} étaient déjà arrivés.

%%K t12-02-2 p
\VUblancAlinea
Comme je voulais envoyer une \VUgras{dépêche (télé}\cc
\VUgras{gramme)} \textsuperscript{(5)}, je me fis indiquer le \VUgras{bureau du télé}\cc
\VUgras{graphe} \textsuperscript{(4)} par un \VUgras{contrôleur} \textsuperscript{(6)}.

%%K t12-03-1 p
\VUblancAlinea
3. --- Avant de passer dans la \VUgras{salle d'attente} \textsuperscript{(7)}\nl
j'allai prendre un \VUgras{billet}\textsuperscript{(9)} d'aller et retour.

%%K t12-04-1 p
\VUblancAlinea
4. --- Je n'attendis pas longtemps au \VUgras{guichet} \textsuperscript{(8)},\nl
car il y avait peu de monde. Un \VUgras{soldat} \textsuperscript{(10)}, qui\nl
partait en permission, me céda obligeamment son\nl
tour, de sorte que j'eus le temps de demander certains\nl
renseignements au \VUgras{chef de gare} \textsuperscript{(13)}, qui causait avec\nl
le \VUgras{commissaire} \textsuperscript{(12)} de surveillance et le \VUgras{gen}\cc
\VUgras{darme} \textsuperscript{(11)} de service.

%%K t12-05-1 p
\VUblancAlinea
5. --- Après avoir acheté un journal à la \VUgras{biblio}\cc
\VUgras{thèque}\textsuperscript{(14)}, je fis peser et enregistrer mes \VUgras{bagages}\textsuperscript{(17)},\nl
qu'un \VUgras{porteur} \textsuperscript{(16)} (facteur) venait d'amener sur un\nl
petit \VUgras{chariot} \textsuperscript{(19)}; l'\VUgras{employé} \textsuperscript{(22)} préposé à la \VUgras{bas}\cc
\VUgras{cule} \textsuperscript{(21)} m'annonça que ma malle pesait 35 kilo
\parplein
\end{VUpage}

% --- Feuillet 72 (folio 70) ------------------------------------------
\begin{VUpage}[72]{70}
%%K t12-05-1 p suite
\VUcontinue
grammes, ce qui me faisait une \VUgras{surcharge} de 5 kilo\cc
grammes, pour laquelle je devais payer un supplé\cc
ment au \VUgras{guichet des bagages} \textsuperscript{(23)}.

%%K t12-06-1 p
\VUblancAlinea
6. --- Comme j'étais en avance, j'eus le loisir d'exa\cc
miner le mouvement des voyageurs dans la gare. Les\nl
uns déposaient ou reprenaient de menus bagages à la\nl
\VUgras{consigne} \textsuperscript{(25)}; d'autres consultaient les \VUgras{horaires} \textsuperscript{(26)}.\nl
Les voyageurs qui arrivaient se pressaient vers la\nl
\VUgras{sortie} \textsuperscript{(32)}, leur \VUgras{valise} \textsuperscript{(24)} à la main. Quelques-uns\nl
s'arrêtaient à la \VUgras{distribution des bagages} \textsuperscript{(27)}\nl
et présentaient leur \VUgras{bulletin}\textsuperscript{(29)} pour retirer leurs\nl
malles, \VUgras{caisses}\textsuperscript{(28)} ou \VUgras{paniers}\textsuperscript{(30)}.

%%K t12-07-1 p
\VUblancAlinea
7. --- Les \VUgras{employés de l'octroi}\textsuperscript{(33)} visitaient soi\cc
gneusement les colis pour voir si l'on avait quelque\nl
chose à déclarer.

%%K t12-08-1 p
\VUblancAlinea
8. --- Certains voyageurs se dirigeaient vers le \VUgras{buf}\cc
\VUgras{fet}\textsuperscript{(34)}, où le \VUgras{garçon}\textsuperscript{(35)} s'empressait de les servir.\nl
Ils étaient forcés de se hâter, car le \VUgras{train}\textsuperscript{(36)}, qui\nl
était un express, ne restait pas longtemps en gare.

%%K t12-09-1 p
\VUblancAlinea
9. --- Je passai ensuite sur le quai de départ et je\nl
cherchai un \VUgras{wagon}\textsuperscript{(37)} direct pour n'être pas obligé de\nl
changer de train pendant le voyage. Je montai dans\nl
une voiture à couloir qui comprend un \VUgras{comparti}\cc
\VUgras{ment}\textsuperscript{(38)} pour fumeurs et un compartiment réservé\nl
aux « dames seules ».

%%K t12-10-1 p
\VUblancAlinea
10. --- Après avoir gravi le \VUgras{marche-pied}\textsuperscript{(40)},\nl
refermé la \VUgras{portière}\textsuperscript{(39)}, relevé le \VUgras{store}\textsuperscript{(41)} et ins\cc
tallé mes menus bagages dans le \VUgras{filet}\textsuperscript{(43)}, je pris\nl
place sur la \VUgras{banquette}\textsuperscript{(42)}. En voyant le \VUgras{lampiste}\textsuperscript{(44)}\nl
allumer les lampes, je songeai que j'avais une nuit à\nl
passer en wagon et j'appelai l'employé chargé de la\nl
location des \VUgras{oreillers}\textsuperscript{(45)} pour lui en demander un.

%%K t12-11-1 p
\VUblancAlinea
11. --- Le conducteur du train vînt vérifier les\nl
billets. Je lui demandai pourquoi nous ne partions pas\nl
puisqu'il était l'heure. Il me répondit que le \VUgras{chef de}\nl
\VUgras{train}\textsuperscript{(46)} était occupé à faire placer des bicyclettes
\parplein
\end{VUpage}

% --- Feuillet 73 (folio 71) ------------------------------------------
\begin{VUpage}[73]{71}
%%K t12-11-1 p suite
\VUcontinue
dans le \VUgras{fourgon des bagages} \textsuperscript{(47)}. De plus, on\nl
n'avait pas encore changé toutes les \VUgras{bouillottes} \textsuperscript{(48)}.

%%K t12-12-1 p
\VUblancAlinea
12. --- Mais nous n'allions pas tarder à partir, car le\nl
\VUgras{chauffeur}\textsuperscript{(50)}, sur son \VUgras{tender}\textsuperscript{(49)}, prenait du charbon\nl
pour activer le feu de la \VUgras{locomotive} \textsuperscript{(54)}, et le \VUgras{méca}\cc
\VUgras{nicien} \textsuperscript{(51)} n'attendait plus que le signal du \VUgras{sous}\cc
\VUgras{chef de gare,} \textsuperscript{(52)} qui était sur le \VUgras{quai} \textsuperscript{(53)}.

%%K t12-13-1 p
\VUblancAlinea
13. --- La \VUgras{chaudière}\textsuperscript{(56)} de la locomotive grondait\nl
sourdement; la \VUgras{cheminée} \textsuperscript{(55)} vomissait des torrents\nl
de fumée mêlée de \VUgras{vapeur}\textsuperscript{(60)}. Un coup de \VUgras{sifflet}\textsuperscript{(59)}\nl
strident retentit et les \VUgras{pistons}\textsuperscript{(58)} se mirent en\nl
mouvement, actionnant les \VUgras{roues}\textsuperscript{(57)} de la lourde\nl
machine.

%%K t12-14-1 p
\VUblancAlinea
14. --- La vitesse du train, filant sur les \VUgras{rails}\textsuperscript{(64)},\nl
s'accéléra; les \VUgras{trottoirs}\textsuperscript{(65)} cessèrent; nous dépas\cc
sâmes les \VUgras{cabinets d'aisances}\textsuperscript{(66)}, ainsi qu'une file\nl
de wagons garés sur une autre \VUgras{voie}\textsuperscript{(63)} : \VUgras{wagons}\cc
\VUgras{lits}\textsuperscript{(67)}, \VUgras{wagon-restaurant}\textsuperscript{(68)}, \VUgras{wagon-poste}\textsuperscript{(69)}.

%%K t12-15-1 p
\VUblancAlinea
15. --- Puis ce fut la \VUgras{gare des marchandises}\textsuperscript{(71)}\nl
qui passa devant nos yeux. Sous les hangars, des\nl
employés chargeaient les \VUgras{marchandises}\textsuperscript{(72)} expé\cc
diées en petite vitesse. Des \VUgras{hommes d'équipe}\textsuperscript{(76)},\nl
près du \VUgras{château-d'eau}\textsuperscript{(73)}, manœuvraient des\nl
\VUgras{wagons à bestiaux}\textsuperscript{(75)} et des \VUgras{wagons plate}\cc
\VUgras{forme}\textsuperscript{(79)}, pour former un \VUgras{train de marchan}\cc
\VUgras{dises}\textsuperscript{(80)}.

%%K t12-16-1 p
\VUblancAlinea
16. --- Nous franchîmes la dernière \VUgras{aiguille}\textsuperscript{(78)},\nl
auprès de laquelle se tenait l'aiguilleur; nous dépas\cc
sâmes le \VUgras{disque}\textsuperscript{(86)} et la gare se perdit dans le\nl
lointain.

%%K t12-17-1 p
\VUblancAlinea
17. --- Bientôt nous nous engageâmes sur le \VUgras{pont}\nl
\VUgras{de fer}\textsuperscript{(74)} qui traverse le \VUgras{fleuve}\textsuperscript{(81)}. Ce pont franchi,\nl
nous longeâmes la rivière, où de puissants \VUgras{remor}\cc
\VUgras{queurs}\textsuperscript{(82)} traînaient de lourds \VUgras{chalands}\textsuperscript{(83)}. Nous\nl
aperçûmes aussi un \VUgras{bateau de voyageurs}\textsuperscript{(84)}, au\nl
moment où celui-ci allait aborder au \VUgras{ponton}\textsuperscript{(85)}.
\end{VUpage}

% --- Feuillet 74 (folio 72) ------------------------------------------
\begin{VUpage}[74]{72}
%%K t12-18-1 p
\VUblancAlinea
18. --- Après avoir brûlé plusieurs stations, nous\nl
traversâmes quelques \VUgras{tunnels}\textsuperscript{(87)} et nous laissâmes\nl
derrière nous d'innombrables \VUgras{maisonnettes}\textsuperscript{(88)} de\nl
\VUgras{gardes-barrières.} Bercé par le cahotement du\nl
train, je m'endormis profondément.

%%K t12-19-1 p
\VUblancAlinea
19. --- Je fus brusquement réveillé par la voix d'un\nl
employé qui criait « Deutsch-Avricourt ! Tout le monde\nl
descend ». C'était la visite de la douane et toutes les\nl
ennuyeuses formalités de la frontière. Il fallut régler\nl
ma montre sur l'\VUgras{horloge}\textsuperscript{(89)} de la gare qui avançait\nl
de 55 minutes; à peine réveillé je distinguais diffici\cc
lement la \VUgras{grande aiguille}\textsuperscript{(90)} de la \VUgras{petite}\textsuperscript{(91)}; le\nl
\VUgras{cadran}\textsuperscript{(92)} lui-même était tout embrouillé par la\nl
brume matinale.

%%K t12-20-1 p
\VUblancAlinea
20. --- Pendant que les douaniers procédaient à leur\nl
visite, je déchiffrai en bâillant les \VUgras{affiches}\textsuperscript{(93)} qui\nl
décorent les murs de la gare. Je déjeunai pour faire\nl
passer le temps, je consultai la carte du \VUgras{réseau}\textsuperscript{(94)}\nl
allemand pour voir quelle distance me séparait encore\nl
de Strasbourg, et je remontai en wagon, réconforté\nl
par l'espoir d'arriver bientôt au but de mon voyage.

\VUsaut{6.0mm}
\VUfilet{22mm}
\end{VUpage}
